Depuis plusieurs ann�es, je me suis fortement investi dans les enseignements et les formations � l'interface entre math�matiques et sciences du vivant. Cette investissement se traduit notamment par des responsabilit�s dans le master {\sl Math�matiques pour les Sciences du Vivant} (MathSV) initialement cr�� par de l'universit� d'Orsay ainsi que dans le programme {\sl Sciences du vivant} du Labex Math�matique Hadamard (LMH) dont l'objectif est notamment de former des math�maticiens de premier rang conscients des principaux concepts et enjeux en biologie et en �cologie,

%%%%%%%%%%%%%%%%%%%%%%%%%%%%%%%%%%%%%%%%%%%%%%%%%%%%%%%%%%%%%%%%%%%%%%%%
\paragraph{Responsabilit�s et activit�s d'enseignement}
% - Conception ou mise en place de projets p�dagogiques.
% - Participation � l'enseignement.
% - Vous pr�ciserez l'intitul� de la formation, le nombre d'heures dispens�es dans la p�riode �valu�e.
\begin{itemize}
 \item (Co-)responsable de la sp�cialisation {\sl Biologie des syst�mes et apprentissage} puis {\sl �cologie et mod�les d'�volution} du master MathSV.
 \item Cours de {\sl Mod�les statistiques pour la g�nomique} au master MathSV  (20h eqTD + r�daction d'un support de cours de 50p). 
% R�daction d'un polycopi� associ�. % \cite{Rob15-MathSV}.
% \item Cours d'{\sl Analyse des donn�es approfondies} (12h eqTD) � l'�cole Nationale de la Statistique et de l'Administration �conomique (ENSAE).
 \item Intervention dans le module {\sl Mod�lisation statistique et analyse de donn�es biologiques} de l'�cole doctorale ABIES (3h de cours, 10h de tutorat).
\end{itemize}

%%%%%%%%%%%%%%%%%%%%%%%%%%%%%%%%%%%%%%%%%%%%%%%%%%%%%%%%%%%%%%%%%%%%%%%%
\paragraph{Responsabilit�s doctorales}
\begin{itemize}
%  \item Membre du conseil de l'�cole doctorale de math�matiques de l'universit� Paris Sud-Orsay.
 \item Correspondant \APT de l'�cole doctorale de math�matiques Jacques Hadamard (EDMH, Paris-Saclay)
\end{itemize}

